\documentclass[11pt]{letter}
\usepackage[margin=1in]{geometry}
\usepackage{hyperref}

\signature{Aaron D. Schroeder\\Social Data Commons\\Biocomplexity Institute and Initiative\\University of Virginia}

\address{Social Data Commons\\Biocomplexity Institute and Initiative\\University of Virginia\\Charlottesville, VA 22904}

\date{\today}

\begin{document}

\begin{letter}{Editorial Office\\Data \& Policy\\Cambridge University Press}

\opening{Dear Editors,}

We submit the enclosed manuscript, ``Licensed childcare accessibility measures for Virginia at the census block group level,'' for consideration as a Data Paper in \textit{Data \& Policy}.

Child care policy in the United States is constrained by a lack of spatially resolved, publicly available data on where licensed supply falls short of demand. Federal and state programs allocate billions of dollars annually in child care subsidies, yet geographic targeting of these investments typically relies on county-level or state-level aggregations that mask substantial within-county variation. This data gap limits the ability of policymakers to identify specific neighborhoods underserved by the licensed child care market and to evaluate whether public investments are reaching the communities with the greatest need.

The dataset described in this manuscript addresses this gap by providing five measures of licensed child care accessibility for the Commonwealth of Virginia at the census block group level for two time points (2021 and 2025). The measures include total licensed capacity, drive time to the nearest provider, and three age-stratified accessibility ratios computed using the three-step floating catchment area method. The dataset covers 5,963 block groups aggregated to 2,198 census tracts, 133 counties, and 35 health districts. It is already integrated into the Virginia Public Health Data dashboard, where it informs public health planning at the county and health district level.

This manuscript is appropriate for \textit{Data \& Policy} because the dataset sits at the intersection of data science and governance. It demonstrates how administrative records, spatial computation, and open data infrastructure can be combined to produce policy-relevant information that would otherwise require bespoke analysis by individual agencies. The block-group-level accessibility measures enable state and local agencies to target child care subsidies toward areas of greatest need, evaluate the spatial equity of facility distributions, and track how the child care supply changed following the COVID-19 pandemic. We believe this aligns with the journal's mission to promote the use of data in an ethical, responsible, and efficient manner for policy.

The manuscript has not been submitted elsewhere and all authors have approved the submission. The data are deposited on Zenodo and freely available for download. All processing code is publicly available at \url{https://github.com/dads2busy/sdc}.

\closing{Sincerely,}

\end{letter}

\end{document}
